% !TeX root = ../../elegantbook-cn.tex
\section{定值定点与非对称韦达}
这部分内容有一定难度，但前面讲过的入题思路和条件翻译在本节中仍然适用。我们依旧按照原有思路进行条件翻译，只是在得到初步结果后，还需要对其进行进一步的“加工”。


\subsection{非对称韦达}
在前面的章节中，我们讲解了通过设直线并运用韦达定理解决解析几何问题的基本思路。此前得到的结果大多可以整理为关于$x_1+x_2$、$x_1x_2$或$y_1+y_2$、$y_1y_2$的对称结构。下面来看一个有所不同的问题：
\begin{tcolorbox}[
  enhanced,
  breakable,
  title={问题引入},
  colback=structurecolor!5,
  colframe=structurecolor,
  colbacktitle=structurecolor,
  coltitle=white,
  fonttitle=\bfseries,
  boxrule=0.5pt,
  arc=1mm,
  left=8pt,
  right=8pt,
  top=6pt,
  bottom=6pt
]
已知
\[
y_1+y_2=-\frac{4n}{n^2+2},\qquad
y_1y_2=\frac{2}{n^2+2},
\]
且$y_1\neq y_2$。求证：$\dfrac{ny_1y_2+y_1}{y_1-y_2}$的值恒为$\dfrac{1}{2}$。
\end{tcolorbox}

这个式子的分子中单独出现了$y_1$，分母中又出现了$y_1-y_2$，二者都不是关于$y_1,y_2$的对称式，因此不能直接利用韦达定理进行整体代入。处理这类问题时，通常有两种思路：一是进行\textbf{和积转换}，二是\textbf{将非对称结构转化为对称结构}。

\medskip
\noindent{\bfseries\color{structurecolor}方法一：和积转换}

韦达定理能够直接给出两根之和$y_1+y_2$与两根之积$y_1y_2$。因此，当已知条件中同时出现$y_1+y_2$、$y_1y_2$和参数$n$时，可以先比较这两个关系式，寻找“和”与“积”之间的联系。
\[
ny_1y_2
=\frac{2n}{n^2+2}
=-\frac{1}{2}(y_1+y_2).
\]
这一步将含参数$n$的乘积$ny_1y_2$转化为了两根之和，从而在两个韦达关系之间建立了联系。将所得关系代入原式，得
\[
\frac{ny_1y_2+y_1}{y_1-y_2}=\frac{-\frac{1}{2}(y_1+y_2)+y_1}{y_1-y_2}=\frac{\frac{1}{2}(y_1-y_2)}{y_1-y_2}=\frac{1}{2}.
\]
\medskip
\noindent{\bfseries\color{structurecolor}方法二：非对称变对称}

题目已经给出了需要证明的定值$\frac{1}{2}$，因此可以将原式与$\frac{1}{2}$作差：
\[
\frac{ny_1y_2+y_1}{y_1-y_2}-\frac{1}{2}=\frac{ny_1y_2+\frac{1}{2}(y_1+y_2)}{y_1-y_2}
\]
作差并通分后，分子由非对称式转化为了对称式。此时可以直接代入韦达定理：
\[
ny_1y_2+\frac{1}{2}(y_1+y_2)
=\frac{2n}{n^2+2}
-\frac{2n}{n^2+2}
=0.
\]
因此，
\[
\frac{ny_1y_2+y_1}{y_1-y_2}=\frac{1}{2}.
\]
这种方法的关键在于：当题目给出了需要证明的定值时，可以将原式与该定值作差，再通过通分、整理，设法将分子转化为关于两根之和与两根之积的对称式。若利用韦达定理能够证明整理后的分子为$0$，原结论便得到了证明。

接下来，再看一个例子。

\begin{tcolorbox}[
  enhanced,
  breakable,
  title={问题引入},
  colback=structurecolor!5,
  colframe=structurecolor,
  colbacktitle=structurecolor,
  coltitle=white,
  fonttitle=\bfseries,
  boxrule=0.5pt,
  arc=1mm,
  left=8pt,
  right=8pt,
  top=6pt,
  bottom=6pt
]
已知
\[
x_1+x_2=\frac{32k^2}{3+4k^2},
\qquad
x_1x_2=\frac{64k^2-12}{3+4k^2}.
\]
证明：$\dfrac{3(x_1-x_2)}{x_1x_2-x_1-4x_2+4}$为定值，并求出该定值。
\end{tcolorbox}

这个问题相对棘手。首先，两根之和与两根之积之间没有明显的转换关系，因此方法一难以直接使用；其次，题目没有给出定值的具体数值，因此方法二也不能直接使用。下面介绍两种可行的思路。

\medskip
\noindent{\bfseries\color{structurecolor}方法三：保留一根，整体代换}

处理非对称韦达问题时，还可以利用恒等式
$x_1=(x_1+x_2)-x_2$
消去$x_1$，使整个式子中只保留$x_2$。同理，也可以消去$x_2$，只保留$x_1$。

具体来说，式子中的$x_1x_2$可以直接利用韦达定理代入；对于非对称的加减部分，则利用$x_1=(x_1+x_2)-x_2$，将所有的$x_1$替换为$x_1+x_2$与$x_2$的组合。随后合并所有不含$x_2$的部分，再观察分子与分母之间是否存在倍数关系。对于本题，有
\[
\begin{aligned}
\frac{3(x_1-x_2)}{x_1x_2-x_1-4x_2+4}
&=\frac{3(x_1+x_2-x_2)-3x_2}
{x_1x_2-(x_1+x_2-x_2)-4x_2+2}\\
&=\frac{3(x_1+x_2)-6x_2}
{x_1x_2-(x_1+x_2)-3x_2+4}\\
&=\frac{\dfrac{96k^2}{3+4k^2}-6x_2}
{\dfrac{64k^2-12}{3+4k^2}
-\dfrac{32k^2}{3+4k^2}-3x_2+4}\\
&=\frac{\dfrac{96k^2}{3+4k^2}-6x_2}
{\dfrac{48k^2}{3+4k^2}-3x_2}
=2.
\end{aligned}
\]


\medskip
\noindent{\bfseries\color{structurecolor}*方法四：根据求根公式确定定值，再作差证明}

可以先确定定值，再按照方法二完成证明。如何确定这个定值呢？由求根公式
$x_1=\dfrac{-b+\sqrt{\Delta}}{2a}$，
$x_2=\dfrac{-b-\sqrt{\Delta}}{2a}$
可以看出，根式部分$\sqrt{\Delta}$不能与其他有理项合并。因此，可以比较原式分子与分母中$\sqrt{\Delta}$项的系数，据此判断定值。

在本题中，分子$3x_1-3x_2$运算后的根式项为$6\sqrt{\Delta}$；分母中含有$\sqrt{\Delta}$的部分来自$-x_1-4x_2$，运算后的根式项为$3\sqrt{\Delta}$。分子与分母中$\sqrt{\Delta}$项的系数之比为$2$，因此可以判断原式的定值为$2$。确定定值后，再将原式与$2$作差，按照方法二将非对称式转化为对称式，最后利用韦达定理完成证明。

\subsection{多个变量的恒等式问题}

当一个含有多个变量或参数的等式对于某个参数的任意取值都成立时，可以将其整理为关于该参数的多项式恒等式，再根据对应项的系数分别相等建立方程组。

\begin{tcolorbox}[
  enhanced,
  breakable,
  title={问题引入},
  colback=structurecolor!5,
  colframe=structurecolor,
  colbacktitle=structurecolor,
  coltitle=white,
  fonttitle=\bfseries,
  boxrule=0.5pt,
  arc=1mm,
  left=8pt,
  right=8pt,
  top=6pt,
  bottom=6pt
]
若
\[
\frac{m^2+(m+1)k^2}{1+2k^2}
\]
对于任意的 \(k\) 都为常数，求 \(m\) 的值。
\end{tcolorbox}

\begin{tcolorbox}[
  enhanced,
  breakable,
  title={问题引入},
  colback=structurecolor!5,
  colframe=structurecolor,
  colbacktitle=structurecolor,
  coltitle=white,
  fonttitle=\bfseries,
  boxrule=0.5pt,
  arc=1mm,
  left=8pt,
  right=8pt,
  top=6pt,
  bottom=6pt
]
若等式
\[
\frac{3kx_0-2y_0-8k^2}{1+2k^2}+x_0^2+y_0^2=0
\]
对于任意的 \(k\) 都成立，求 \(x_0,y_0\) 的值。
\end{tcolorbox}

\subsection{定值}
接下来正式进入定值问题。解决这类问题需要把握一个关键环节：如果最后需要证明的表达式是非对称的，不如先选取便于处理的特殊情形，例如令直线斜率不存在、斜率为$0$，或者利用图形的对称性，初步判断定值的具体数值，然后使用非对称韦达的方法二完成证明。




\begin{example} % 定值定点，特殊法先得答案，非对称韦达
（2025朝阳一模）已知椭圆 $E:\frac{x^{2}}{a^{2}}+\frac{y^{2}}{b^{2}}=1\ (a>b>0)$ 的右焦点为 $F(1,0)$，离心率为 $\frac{1}{2}$.\\
(1) 求椭圆 $E$ 的方程；\\
(2) 过点 $M(4,0)$ 作直线 $l$ 与椭圆 $E$ 交于不同的两点 $A,B$. 设 $C\left(\frac{5}{2},\frac{3}{2}\right)$，直线 $BC$ 与直线 $x=1$ 交于点 $N$. 求证：直线 $AN$ 的斜率为定值.
\end{example}
\begin{example} % 直线引入，定值定点与非对称韦达
（2022朝阳一模）已知椭圆 $C:\frac{x^2}{a^2}+\frac{y^2}{b^2}=1\ (a>b>0)$ 的一个焦点为 $F(1,0)$，且过点 $\left(1,\frac{3}{2}\right)$.\\
（1）求椭圆 $C$ 的方程和离心率；\\
（2）过点 $P(4,0)$ 且与 $x$ 轴不重合的直线 $l$ 与椭圆 $C$ 交于 $A,B$ 两点，与直线 $x=1$ 交于点 $Q$，点 $M$ 满足 $MP\perp x$ 轴，$MB\parallel x$ 轴，试求直线 $MA$ 的斜率与直线 $MQ$ 的斜率的比值.
\end{example}

\begin{example} 
（2020北京高考）已知椭圆 $C:\frac{x^2}{a^2}+\frac{y^2}{b^2}=1$ 过点 $A(-2,-1)$，且 $a=2b$.\\
（I）求椭圆 $C$ 的方程；\\
（II）过点 $B(-4,0)$ 的直线交椭圆 $C$ 于点 $M,N$，直线 $MA,NA$ 分别交直线 $x=-4$ 于点 $P,Q$. 求 $\frac{\lvert PB\rvert}{\lvert BQ\rvert}$ 的值.
\end{example}
\begin{example} % 直线引入，定值定点（两个未知数的恒等式）
（2021丰台二模）已知椭圆 $C:\frac{x^2}{3}+y^2=1$，过点 $(-1,0)$ 的直线 $l$ 交椭圆 $C$ 于点 $A,B$。\\
（1）当直线 $l$ 与 $x$ 轴垂直时，求 $\lvert AB\rvert$；\\
（2）在 $x$ 轴上是否存在定点 $P$，使 $\overrightarrow{PA}\cdot\overrightarrow{PB}$ 为定值？若存在，求点 $P$ 的坐标及 $\overrightarrow{PA}\cdot\overrightarrow{PB}$ 的值；若不存在，说明理由。
\end{example}







\subsection{定点}

\medskip
\noindent{\bfseries\color{structurecolor}直线过定点}


先观察下面三个简单的直线族：
\[
y=kx+1,\qquad
y=kx-k+3,\qquad
x=ny+2n-1.
\]
虽然直线的位置会随参数变化，但每个直线族中都隐藏着一个不随参数变化的点。将方程适当整理，得到
\[
y=kx+1,\qquad
y-3=k(x-1),\qquad
x+1=n(y+2).
\]
由此可以直接看出：第一组直线恒过定点$(0,1)$，第二组直线恒过定点$(1,3)$，第三组直线恒过定点$(-1,-2)$。

\textbf{方法一：寻找斜率与截距的关系}

一般地，将直线写成
\[
y=kx+m
\]
的形式后，可以寻找斜率$k$与截距$m$之间的关系，再将直线方程整理为
\[
y-y_0=k(x-x_0).
\]
其中不含参数的点$(x_0,y_0)$，就是直线恒过的定点。

例如，当$k=m$时，
\[
y=kx+m=kx+k=k(x+1),
\]
因此，无论$k$取何值，直线都恒过定点$(-1,0)$。

\textbf{方法二：特殊情形确定方向}


如果直线方程较为复杂，例如
\[
y=\frac{y_1-y_2}{1-x_2}(x-1)+y_1,
\]
往往很难直接得到斜率与截距之间的关系。此时可以先选取便于计算的特殊情形，初步判断定点的位置，或者判断定点位于哪一条直线上。

例如，若通过特殊情形判断定点在$x$轴上，则可设该点为$P(x_0,0)$。接下来在一般情形下令$y=0$，再证明所得的$x_0$与参数无关即可。这样，原来的定点问题就被转化为了定值问题。如果整理过程中出现单独的$x_1$、$x_2$或两根之差等非对称结构，则继续使用前面介绍的非对称韦达方法进行转化。


\begin{example} % 直线引入，
（2026丰台二模）已知椭圆 $E:\frac{x^2}{a^2}+\frac{y^2}{b^2}=1\ (a>b>0)$ 的离心率为 $\frac{\sqrt{3}}{2}$，以椭圆 $E$ 的焦点和短轴顶点为顶点的四边形是边长为 $2$ 的菱形。\\
（1）求椭圆 $E$ 的标准方程；\\
（2）已知 $A$ 为椭圆 $E$ 的左顶点，$M,N$ 为椭圆 $E$ 上两个不同的动点（均不与点 $A$ 重合），且直线 $AM$ 与直线 $AN$ 的斜率之积为 $\frac{5}{4}$。求证：直线 $MN$ 过定点。
\end{example}

\begin{example} % 直线引入，菱形的翻译
    % 平行四边形例题
（2024海淀二模）已知椭圆 $E$ 的焦点在 $x$ 轴上，中心在坐标原点。以 $E$ 的一个顶点和两个焦点为顶点的三角形是等边三角形，且其周长为 $6\sqrt{2}$。\\
（1）求椭圆 $E$ 的方程；\\
（2）设过点 $M(2,0)$ 的直线 $l$（不与坐标轴垂直）与椭圆 $E$ 交于不同的两点 $A,C$，与直线 $x=16$ 交于点 $P$。点 $B$ 在 $y$ 轴上，$D$ 为坐标平面内的一点，四边形 $ABCD$ 是菱形。求证：直线 $PD$ 过定点。
\end{example}
\begin{example} % 入题思路（我搞复杂了），直线过顶点
% 入题思路例题1.14，定点例题
（2022朝阳二模）已知椭圆 $C:\frac{x^2}{a^2}+\frac{y^2}{b^2}=1\ (a>b>0)$ 的一个顶点为 $P(0,1)$，离心率为 $\frac{\sqrt{2}}{2}$。\\
（1）求椭圆 $C$ 的方程；\\
（2）过点 $P$ 作斜率为 $k_1$ 的直线 $l_1$，交椭圆 $C$ 于另一点 $A$；过点 $P$ 作斜率为 $k_2\ (k_2\neq k_1)$ 的直线 $l_2$，交椭圆 $C$ 于另一点 $B$。若 $k_1k_2=1$，求证：直线 $AB$ 经过定点。
\end{example}
\begin{example} % 直线引入，三点共线，定值定点（不需要什么技巧）
（2022海淀一模）已知椭圆 $C:\frac{x^2}{a^2}+\frac{y^2}{b^2}=1\ (a>b>0)$ 的下顶点 $A$ 和右顶点 $B$ 都在直线 $l_1:y=\frac{1}{2}(x-2)$ 上.\\
（1）求椭圆方程及其离心率；\\
（2）不经过点 $B$ 的直线 $l_2:y=kx+m$ 交椭圆 $C$ 于两点 $P,Q$，过点 $P$ 作 $x$ 轴的垂线交 $l_1$ 于点 $D$，点 $P$ 关于点 $D$ 的对称点为 $E$. 若 $E,B,Q$ 三点共线，求证：直线 $l_2$ 经过定点.
\end{example}
\begin{example} % 直线，定点，非对称韦达
（2021西城一模）已知椭圆 $C:\frac{x^2}{a^2}+\frac{y^2}{3}=1\ (a>0)$ 的焦点在 $x$ 轴上，且经过点 $E\left(1,\frac{3}{2}\right)$，左顶点为 $D$，右焦点为 $F$.\\
（1）求椭圆 $C$ 的离心率和 $\triangle DEF$ 的面积；\\
（2）已知直线 $y=kx+1$ 与椭圆 $C$ 交于 $A,B$ 两点，过点 $B$ 作直线 $y=t\ (t>\sqrt{3})$ 的垂线，垂足为 $G$. 判断是否存在常数 $t$，使得直线 $AG$ 经过 $y$ 轴上的定点？若存在，求 $t$ 的值；若不存在，请说明理由.
\end{example}


\medskip
\noindent{\bfseries\color{structurecolor}以线段为直径的圆恒过定点}

对于以线段为直径的圆恒过定点问题，通常有以下两种思路。

第一种思路是先通过参数的特殊取值猜出定点，再证明该点始终在圆上。证明时，需要将“点在以线段为直径的圆上”翻译为：该点与直径两端点的连线互相垂直。

第二种思路是直接设定点为\(P(x_0,y_0)\)。如果根据题目条件能够判断定点在某条特殊直线上，还可以进一步简化设点。例如，若定点在\(y\)轴上，则可设\(P(0,y_0)\)。随后，将“点\(P\)在以线段为直径的圆上”翻译为相应的等式，整理出关于参数的恒等式，再利用前面介绍的多个变量的恒等式问题求出\(x_0,y_0\)。









\begin{example} % 直线引入，定值定点，长度转圆
    % 练习题
（2024丰台一模）已知椭圆 $E:\frac{x^2}{a^2}+\frac{y^2}{b^2}=1\ (a>b>0)$ 的焦距为 $4\sqrt{2}$，以椭圆 $E$ 的四个顶点为顶点的四边形的周长为 $16$.\\
（1）求椭圆 $E$ 的标准方程；\\
（2）过点 $S(0,1)$ 的直线 $l$ 交椭圆 $E$ 于 $P,Q$ 两点，线段 $PQ$ 的中点为 $M$. 是否存在定点 $D$，使得 $\frac{\lvert DM\rvert}{\lvert PQ\rvert}=\frac{1}{2}$？若存在，求出 $D$ 的坐标；若不存在，请说明理由.
\end{example}









% \section{怎么条件转换最优？}

% \begin{exercise} % 中点与平行条件转换为斜率关系
% （2022东城期末）已知椭圆 $C:\frac{x^2}{a^2}+\frac{y^2}{b^2}=1$ 过点 $A(-\sqrt{3},0)$，其右焦点为 $F(1,0)$.\\
% （1）求椭圆 $C$ 的方程；\\
% （2）设 $P$ 为椭圆 $C$ 上一动点（不在 $x$ 轴上），$M$ 为 $AP$ 中点，过原点 $O$ 作 $AP$ 的平行线，与直线 $x=3$ 交于点 $Q$. 问：直线 $OM$ 与 $FQ$ 斜率的乘积是否为定值？若为定值，求出该值；若不为定值，请说明理由.
% \end{exercise}
